%%
%% AREST: Applicative REpresentational State Transfer
%% Prepared for submission to Communications of the ACM
%% Using acmsmall template per CACM Author Guidelines
%%
\documentclass[acmsmall,screen,review]{acmart}

\AtBeginDocument{%
  \providecommand\BibTeX{{Bib\TeX}}}

%% Rights management placeholder (to be completed upon acceptance)
\setcopyright{acmlicensed}
\copyrightyear{2026}
\acmYear{2026}
\acmDOI{XXXXXXX.XXXXXXX}
\acmJournal{JACM}
\acmVolume{0}
\acmNumber{0}
\acmArticle{0}
\acmMonth{0}

\usepackage{listings}
\usepackage{enumitem}
\lstset{
  basicstyle=\small\ttfamily,
  breaklines=true,
  frame=single,
  xleftmargin=1em,
  xrightmargin=1em,
}

%%
%% Title
%%
\title{Compiling Facts into Applications}

%%
%% Author
%%
\author{Samuel Lippert}
\email{samuel@driv.ly}
\orcid{0009-0003-9048-9857}
\affiliation{%
  \institution{Drivly Inc}
  \city{Shakopee}
  \state{Minnesota}
  \country{USA}
}

%%
%% Abstract
%%
\begin{abstract}
We define AREST, a compositional architecture unifying Backus's Formal Functional Programming and Applicative State Transition systems~\cite{backus78} with Fielding's Representational State Transfer~\cite{fielding00}, using Halpin's ORM\,2~\cite{halpin08} as the surface syntax. Domain knowledge is expressed as FFP objects in FORML\,2~\cite{halpin06}; the representation function~$\rho$ maps each object to its executable; the system state is Backus's AST cell sequence; the persistent data is Codd's named set~\cite{codd70}. We prove eight properties: grammar unambiguity, specification equivalence, completeness of state transfer, HATEOAS as projection, derivability, query adequacy, derivation completeness, and scope.
\end{abstract}

%%
%% CCS Concepts
%%
\begin{CCSXML}
<ccs2012>
 <concept>
  <concept_id>10011007.10011006.10011039</concept_id>
  <concept_desc>Software and its engineering~Formal language definitions</concept_desc>
  <concept_significance>500</concept_significance>
 </concept>
 <concept>
  <concept_id>10011007.10011006.10011050.10011017</concept_id>
  <concept_desc>Software and its engineering~Domain specific languages</concept_desc>
  <concept_significance>500</concept_significance>
 </concept>
 <concept>
  <concept_id>10003752.10003790.10003795</concept_id>
  <concept_desc>Theory of computation~Constraint and logic programming</concept_desc>
  <concept_significance>300</concept_significance>
 </concept>
 <concept>
  <concept_id>10002951.10003260.10003282</concept_id>
  <concept_desc>Information systems~RESTful web services</concept_desc>
  <concept_significance>300</concept_significance>
 </concept>
</ccs2012>
\end{CCSXML}

\ccsdesc[500]{Software and its engineering~Formal language definitions}
\ccsdesc[500]{Software and its engineering~Domain specific languages}
\ccsdesc[300]{Theory of computation~Constraint and logic programming}
\ccsdesc[300]{Information systems~RESTful web services}

%%
%% Keywords
%%
\keywords{applicative state transfer, REST, formal functional programming, object-role modeling, state machines, HATEOAS, first-order logic, specification equivalence}

%%
%% Build document
%%
\begin{document}

\maketitle

\section{Definitions}

\begin{definition}[AREST]\label{def:arest}
An AREST system is a tuple $(O, \rho, D, P, \mathit{sub})$ where:
\begin{itemize}[nosep]
\item $O$ is the set of FFP objects~\cite[Sec.\,13.2]{backus78}
\item $\rho: O \to (O \to O)$ is the representation function~\cite[Sec.\,13.3.2]{backus78}
\item $D$ is the AST state, a sequence of cells~\cite[Sec.\,14.3]{backus78}
\item $P = \uparrow\!\mathit{FILE}\!:\!D$ is the population, a named set of relations~\cite[Sec.\,1.6]{codd70} whose elements are ORM\,2 elementary facts~\cite[Ch.\,2]{halpin08}
\item $\mathit{sub}$ is the subsystem function dispatching inputs to state transitions, producing output as REST representations~\cite{fielding00} with HATEOAS links
\end{itemize}
\end{definition}

Three terms carry overloaded abbreviations. \emph{AST} follows Backus's usage (Applicative State Transition, 1978), not the compiler term. \emph{ORM} follows Halpin's usage (Object-Role Modeling, 2008), not the persistence pattern. \emph{FP} follows Backus's usage (Function-level Programming), not the broader tradition.

The primitive atoms of $O$ include Backus's structural functions~\cite[Sec.\,11.2.3]{backus78}. Sequences represent functional forms via metacomposition~\cite[Sec.\,13.3.2]{backus78}:
\begin{align}
(\rho\langle x_1, \ldots, x_n\rangle)\!:\!y = (\rho\,x_1)\!:\!\langle\langle x_1, \ldots, x_n\rangle, y\rangle \label{eq:metacomp}
\end{align}
The first element determines which functional form; the remaining elements are its parameters. In AREST, a FORML\,2 reading is an object; $\rho$ extracts the constraint or derivation function it represents.

\section{Operations}

\subsection{Domain Primitives}

Each combining form~\cite[Sec.\,11.2.4]{backus78} is an FFP object with a domain interpretation:

\begin{table}[h]
\caption{Backus combining forms as FFP objects with AREST domain interpretations.}
\label{tab:forms}
\centering
\small
\begin{tabular}{@{}lll@{}}
\toprule
\textbf{FFP Object} & \textbf{Form} & \textbf{Domain} \\
\midrule
$\langle\mathit{CONS}, s_1, \ldots, s_n\rangle$ & Construction & Graph Schema \\
$s_i$ & Selector & Role \\
$\langle\mathit{COMP}, f, g\rangle$ & Composition & Derivation \\
$\langle\mathit{COND}, p, f, g\rangle$ & Condition & Constraint \\
$\langle\alpha, f\rangle$ & Apply-to-all & Traversal \\
$\langle/, f\rangle$ & Insert & Aggregation \\
$\langle\mathit{bu}, f, x\rangle$ & Binary-to-unary & Query \\
$\bar{x}$ & Constant & Literal \\
\bottomrule
\end{tabular}
\end{table}

A graph schema $\langle\mathit{CONS}, s_1, \ldots, s_n\rangle$ applied to all roles is a fact; partially applied via $\mathit{bu}$ it is a query.

\subsection{Relational Algebra}

Codd's adequate $\theta_1$~\cite[Sec.\,2.2]{codd70} as FFP definitions:
\begin{align}
\pi_L(R) &\equiv \alpha\,[s_{i_1}, \ldots, s_{i_k}] : R \label{eq:proj} \\
R \mathbin{*} S &\equiv \mathit{Filter}(\mathit{eq} \circ [s_{\mathit{sh}}]) \circ \mathit{distl} \label{eq:join} \\
R_{L|_M}\!S &\equiv \mathit{Filter}(p) : R \label{eq:restrict} \\
R \cdot S &\equiv \pi_{1s}(R \mathbin{*} S) \label{eq:rcomp} \\
\gamma(R) &\equiv \mathit{Filter}(\mathit{eq} \circ [s_1, s_n]) : R \label{eq:tie}
\end{align}

\subsection{State and Cells}

The state $D$ is a sequence of cells $\langle\mathit{CELL}, n, c\rangle$~\cite[Sec.\,14.3]{backus78}:
\begin{align}
\uparrow\!n\!:\!D &\to c \qquad \text{(fetch)} \label{eq:fetch} \\
\downarrow\!n\!:\!\langle x, D\rangle &\to D' \quad\; \text{(store)} \label{eq:store}
\end{align}
$\mathit{FILE}$ is a cell whose contents is~$P$. Each entity is a cell. Cells may contain sub-stores~\cite[Sec.\,14.7]{backus78}; partitioning is a cell containing a subset of~$P$.

\subsection{Subsystem Dispatch}

Following~\cite[Sec.\,14.4.2]{backus78}:
\begin{align}
&\mathit{sub} \equiv \nonumber \\
&\; \mathit{is\text{-}chg} \circ \uparrow\!\mathit{K} \to [\mathit{rpt}, \mathit{apl}] \circ [\uparrow\!\mathit{I}, \mathit{defs}]; \nonumber \\
&\; \mathit{is\text{-}qry} \circ \uparrow\!\mathit{K} \to [\mathit{qry} \circ [\uparrow\!\mathit{I}, \uparrow\!\mathit{F}], \mathit{defs}]; \nonumber \\
&\; \mathit{is\text{-}cmd} \circ \uparrow\!\mathit{K} \to \mathit{chk} \circ \mathit{apl} \circ [\uparrow\!\mathit{I}, \mathit{defs}]; \nonumber \\
&\; \mathit{is\text{-}upd} \circ \uparrow\!\mathit{K} \to \nonumber \\
&\;\quad [\mathit{rpt}, \downarrow\!\mathit{F} \circ [\mathit{upd}, \mathit{defs}]] \circ [\uparrow\!\mathit{I}, \uparrow\!\mathit{F}]; \nonumber \\
&\; [\mathit{err} \circ [\uparrow\!\mathit{K}, \uparrow\!\mathit{I}], \mathit{defs}] \label{eq:sub}
\end{align}
where $K = \mathit{KEY}$, $I = \mathit{INPUT}$, $F = \mathit{FILE}$.

\subsection{Commands}

Each interaction is an AST transition~\cite[Sec.\,14.3.1]{backus78}:
\begin{align}
\mu(\mathit{SYSTEM}\!:\!x) &= \langle o, D' \rangle \label{eq:transition}
\end{align}

A \textsc{Create} command:
\begin{align}
\mathit{create} &= \mathit{emit} \circ \mathit{validate} \circ \mathit{derive} \circ \mathit{resolve} \label{eq:create}
\end{align}

A state machine is a fold~\cite[Sec.\,11.2.4]{backus78}:
\begin{align}
\mathit{machine} &= /\mathit{transition} : \langle e_1, \ldots, e_n \rangle \label{eq:sm} \\
\mathit{transition} &= (p \to \bar{s'} ;\; \mathit{id}) \label{eq:trans}
\end{align}

\section{Example}

An order domain is declared in FORML\,2:

\begin{lstlisting}
Order(.OrderId) is an entity type.
Customer(.Name) is an entity type.
Order was placed by Customer.
  Each Order was placed by exactly one
  Customer.
Status 'Draft' is initial in State Machine
  Definition 'Order'.
Transition 'place' is from Status 'Draft'.
  Transition 'place' is to Status 'Placed'.
\end{lstlisting}

The constraint reading is an FFP object; $\rho$ yields:
\begin{align*}
(\rho\,\mathit{reading}) = \langle\mathit{COND},\; &(\mathit{count} \leq 1) \circ \mathit{Filter}(\mathit{eq} \circ [s_1, \bar{o}]),\\
&\bar{\phi},\; \mathit{violation}\rangle : P_{\mathit{placedBy}}
\end{align*}

\texttt{POST /orders \{"customer":"acme"\}} dispatches as $\mathit{is\text{-}cmd}$ in~(\ref{eq:sub}). The engine applies $\mathit{create}$~(\ref{eq:create}):
\begin{align*}
P' = P \cup \{&\mathit{Order}(\texttt{ord\text{-}1}),\;
\mathit{placedBy}(\texttt{ord\text{-}1}, \texttt{acme}),\\
&\mathit{SM}(\texttt{ord\text{-}1}, \texttt{Draft})\}
\end{align*}

Response (Theorem~\ref{thm:hateoas}: links from $P'$):
\begin{lstlisting}
{ "id": "ord-1", "customer": "acme",
  "status": "Draft",
  "_links": {
    "place": {
      "href": "/orders/ord-1/transition",
      "method": "POST" } } }
\end{lstlisting}

Following ``place'' applies~(\ref{eq:sm}). The new representation contains $\{\mathit{ship}\}$, not $\{\mathit{place}\}$. If a second command assigns a different Customer to the same Order, the UC fires and the command is rejected. The violation is returned as the original reading.

\section{Execution Semantics}

\subsection{Evaluation as $\beta$-Reduction}

FFP objects denote $\lambda$-terms via Backus's correspondence~\cite[Sec.\,12.9]{backus78}. The representation function $\rho$ yields the executable term. Execution is applicative-order (innermost-first) $\beta$-reduction to normal form:

$\mu$ evaluates an expression under definitions $D$ by repeatedly replacing its innermost application $(x\!:\!y)$ with $(\rho\,x)\!:\!y$, where $\rho$ resolves $x$ to a $\lambda$-term under environment $D$ and primitives over $P$~\cite[Sec.\,13.3.1]{backus78}. Each step is a $\beta$-reduction: $\rho$ yields the $\lambda$-term; application reduces it. The function $\mathit{apply}(\mathit{func}, x, D)$ in the implementation is this reduction step.

All evaluation in AREST is over finite $P$ with monotonic derivations; reduction terminates for all requests.

\subsection{Evaluation Strategy}

Three constraints govern execution:
\begin{enumerate}[nosep]
\item \textbf{Reduction order:} applicative (call-by-value). All arguments are reduced before the function is applied.
\item \textbf{Scope:} evaluation is confined to the request-local partition of $D$. RMAP~\cite[Ch.\,17]{halpin08} determines the partition from the schema's uniqueness constraints.
\item \textbf{Derivation:} forward chaining, monotonic, evaluated to the least fixed point per request. Since $\mathit{derive}$ only adds facts over finite $P_{\mathit{local}}$, the fixed point exists (Knaster-Tarski) and is reached in $\leq |P_{\mathit{local,max}}|$ iterations.
\end{enumerate}
The completeness result (Theorem~\ref{thm:complete}) assumes evaluation to the least fixed point under this strategy.

\subsection{Boundary Conditions}

The system operates over three categories:
\begin{enumerate}[nosep]
\item \textbf{Facts} $(P)$: ORM\,2 elementary facts, persisted in cells.
\item \textbf{Computations} $(O$ via $\rho)$: total, deterministic functions over $P$.
\item \textbf{Inputs} $(I)$: request payload, timestamp, external signals --- part of the operand $x$ in $\mu(\mathit{SYSTEM}\!:\!x)$~\cite[Sec.\,14.3.1]{backus78}.
\end{enumerate}

Time, randomness, and external effects are classified as follows:
\begin{itemize}[nosep]
\item \textbf{Time}: an input fact. ``Now'' is part of $x$, not hidden state. Time-dependent queries are constraints on timestamp facts in $P$.
\item \textbf{Effects}: the output $o$ in $\langle o, D'\rangle$. The engine computes directives (send email, charge payment); the runtime executes them. Outcomes are facts in $D'$.
\item \textbf{Non-determinism}: identity generation is a primitive function producing a fresh atom, or the client provides the ID in $I$.
\item \textbf{Partial information}: the world assumption per noun~\cite[Ch.\,7]{halpin08} governs incompleteness. CWA: not stated $=$ false. OWA: not stated $=$ unknown.
\item \textbf{Identity}: the reference scheme~\cite[Ch.\,5]{halpin08} determines identity before facts are asserted. $\mathit{resolve}$ in~(\ref{eq:create}) applies the selector.
\end{itemize}

AREST governs the derivation of representations from $(P, I)$. Theorem~\ref{thm:scope} holds: every value derivable from $(P, I)$ by a computable function is in scope; values outside are non-factual or non-computable.

\section{Properties}

\begin{theorem}[Grammar Unambiguity]\label{thm:grammar}
Let $N$ be a declared set of object type names and $F$ a declared set of predicate readings, such that no element of $N$ contains a formal item of the grammar as a substring. Then each FORML\,2 sentence over $(N, F)$ has exactly one parse.
\end{theorem}

\begin{proof}
A FORML\,2 sentence has the structure:
\begin{align*}
[\mathit{modal}] \;\; [\mathit{quant}_1] \;\; n_1 \;\; r \;\; \mathit{quant}_2 \;\; n_2 \;\; [\mathit{qual}]
\end{align*}
where $\mathit{modal} \in \{$``It is obligatory that'', ``It is forbidden that'', ``It is permitted that'', $\varepsilon\}$; $\mathit{quant}_1 \in \{$``Each'', ``For each'', ``the same'', $\varepsilon\}$; $n_1, n_2 \in N$; $r \in F$; $\mathit{quant}_2 \in \{$``at most one'', ``some'', ``exactly one'', ``at least $k$ and at most $m$'', ``more than one'', ``no'', $\varepsilon\}$; $\mathit{qual} \in \{$``who $\ldots$'', ``that $\ldots$'', $\varepsilon\}$.

\medskip\noindent
(1)~The modal prefixes are pairwise non-overlapping strings, so $\mathit{modal}$ is determined by the sentence prefix.

(2)~After $\mathit{modal}$, the quantifier $\mathit{quant}_1$ is determined: ``Each'' vs.\ ``For each'' vs.\ ``the same'' vs.\ a name from $N$ are distinguishable because no $n \in N$ contains a formal item as a substring (by hypothesis).

(3)~$n_1$ is determined by longest-first matching against $N$. Since names do not contain formal items, the boundary between $n_1$ and $r$ is unambiguous.

(4)~$r$ is the maximal token sequence between $n_1$ and $\mathit{quant}_2$. The quantifier phrases in $\mathit{quant}_2$ are pairwise distinct strings: \{``at most one'', ``some'', ``exactly one'', ``at least'', ``more than one'', ``no''\}. None is a prefix of another in this position. Since no $n \in N$ contains these phrases, the boundary between $r$ and $\mathit{quant}_2$ is unambiguous.

(5)~$n_2$ is determined by longest-first matching against $N$ after $\mathit{quant}_2$.

(6)~$\mathit{qual}$ is determined by ``who'' or ``that'' (personal vs.\ impersonal pronouns per~\cite[Sec.\,1.5]{halpin06b}).

Each structural position admits exactly one parse. The constraint type is determined by the pair $(\mathit{quant}_1, \mathit{quant}_2)$, which maps injectively to constraint kinds: (Each, at most one)~$\to$~UC; (Each, some)~$\to$~MC; (Each, exactly one)~$\to$~UC$+$MC; (Each, at least $k$\ldots)~$\to$~FC; (the same, more than one)~$\to$~UC$^{-}$. No two pairs map to the same kind.
\end{proof}

\begin{theorem}[Specification Equivalence]\label{thm:spec}
Let $\mathit{parse}: R \to \Phi$ and $\mathit{compile}: \Phi \to O$ be injective, with $R$ the set of FORML\,2 readings, $\Phi$ the set of FOL formulas, and $O$ the set of FFP objects. Then:
\begin{align}
\mathit{parse}^{-1} \circ\, \mathit{compile}^{-1} \circ\, \mathit{compile} \circ \mathit{parse} &= \mathit{id}_R \label{eq:spec}
\end{align}
and $\rho(\mathit{compile}(\mathit{parse}(r)))$ is the executable for every $r \in R$.
\end{theorem}

\begin{proof}
$\mathit{parse}$ is injective by Theorem~\ref{thm:grammar}: each sentence has one parse, so distinct sentences yield distinct formulas. $\mathit{compile}$ is injective: the $(\mathit{quant}_1, \mathit{quant}_2)$ pair determines the constraint kind (Theorem~\ref{thm:grammar}), which determines the FFP object structure (Condition for UC, null-test for MC, count for FC). Both have left inverses; composing: $\mathit{parse}^{-1} \circ \mathit{compile}^{-1} \circ \mathit{compile} \circ \mathit{parse} = \mathit{id}_R$. Executability: $\rho$ maps the compiled object to a function by~\cite[Sec.\,13.4]{backus78}.
\end{proof}

\begin{theorem}[Completeness of State Transfer]\label{thm:complete}
$\mathit{create}$~(\ref{eq:create}) produces a complete $\langle o, D'\rangle$ containing entity, state machine instance, derived facts, violations, and HATEOAS links.
\end{theorem}

\begin{proof}
Let $P = \uparrow\!\mathit{FILE}\!:\!D$. (1)~$\mathit{resolve}$ inserts entity facts: $P' = P \cup E$. (2)~$\mathit{derive}$ is monotonic over finite~$P$; by Knaster-Tarski, a unique least fixed point $P''$ exists. (3)~$\mathit{validate}$ evaluates all constraint objects in $D$ against $P''$: $V = \{(\rho\,c)\!:\!P'' \mid c \in D_{\mathit{constraints}}\}$. (4)~$\mathit{emit}$ produces $\langle o, D[\mathit{FILE} \mapsto P'']\rangle$ if no alethic violation, else $D' = D$. The state does not change during evaluation~\cite[Sec.\,14.1]{backus78}.
\end{proof}

\begin{theorem}[HATEOAS as Projection]\label{thm:hateoas}
Let $T = \{(s_i, s_j, e_k)\} \subseteq P$ be transition facts and $s$ the current status. Then:
\begin{align}
\mathit{links}(s) &= \pi_{\mathit{event}}\!\left(\mathit{Filter}(p) : T\right) \label{eq:hateoas}
\end{align}
where $p(t) = (s_{\mathit{from}}(t) = s) \lor \mathit{anc}(s_{\mathit{from}}(t), s)$.
\end{theorem}

\begin{proof}
$T \subseteq P$. $\mathit{Filter}$ is restriction~(\ref{eq:restrict}); $\pi$ is projection~(\ref{eq:proj}). Both are $\theta_1$ operations on the named set. The representation renders $P$; $\mathit{links}(s)$ is derived from $P$ by $\theta_1$ operations and therefore appears in the representation without external computation.
\end{proof}

\begin{theorem}[Derivability]\label{thm:derive}
For every domain value $v$ in an AREST representation:
\begin{align}
\exists f \in O.\; v = (\rho\,f)\!:\!P \label{eq:derive}
\end{align}
\end{theorem}

\begin{proof}
Domain values partition into: (1)~entity attributes (facts in $P$; $v = (\rho\,s_i)\!:\!\mathit{fact}$), (2)~derived facts ($v = (\rho\,r)\!:\!P$ for derivation rule $r$), (3)~state machine status (derived fact), (4)~links ($v = (\rho\,\langle\mathit{COMP}, \pi, \mathit{Filter}(p)\rangle)\!:\!T$ by Theorem~\ref{thm:hateoas}), (5)~violations ($v = (\rho\,c)\!:\!P$ for constraint $c$). Each case yields~(\ref{eq:derive}).
\end{proof}

\begin{theorem}[Query Adequacy]\label{thm:adequate}
Every query on $P$ is expressible using $\theta_1$, each element of which is an FFP definition~(\ref{eq:proj})--(\ref{eq:tie}).
\end{theorem}

\begin{proof}
$P$ is a named set of relations in normal form (ORM\,2 elementary facts are 5NF by construction~\cite[Ch.\,12]{halpin08}). $\theta_1$ is adequate for the named set~\cite[Sec.\,2.2]{codd70}. Each element of $\theta_1$ is an FFP definition by~(\ref{eq:proj})--(\ref{eq:tie}).
\end{proof}

\begin{theorem}[Derivation Completeness]\label{thm:deriv-complete}
Every $\theta_1$-derivable relation has a FORML\,2 reading. Forward chaining evaluates the corresponding FFP objects to the least fixed point.
\end{theorem}

\begin{proof}
Each $\theta_1$ operation has a FORML\,2 verbalization: projection verbalizes as role selection (``$X$ has $Y$''); join as shared-role composition (``$X$ who $R_1$ $Y$ and $R_2$ $Z$''); restriction as constraint predicate (``$X$ where $Y$ exceeds $Z$''); tie as self-reference (``$X$ is ancestor of $X$''). By Theorem~\ref{thm:grammar}, each verbalization has a unique parse, so each $\theta_1$ composition has a unique reading. Forward chaining applies derivation rules (FFP objects) to $P$ monotonically; by Knaster-Tarski the least fixed point exists and is reached in $|P_{\max}|$ steps.
\end{proof}

\begin{theorem}[Scope]\label{thm:scope}
The architecture covers all values that are both factual and computable. Values outside are either non-factual or non-computable.
\end{theorem}

\begin{proof}
Any factual proposition (a statement with a truth value) is expressible as a tuple in a relation~\cite[Sec.\,1.3]{codd70}, hence a fact in $P$. Any computable function (an effective method) is expressible as a $\lambda$-term~\cite{church36}; Backus shows the equivalence of FP combining forms to $\lambda$-calculus combinators~\cite[Sec.\,12.9]{backus78}, and FFP extends FP with metacomposition. Therefore every factual-and-computable value is $(\rho\,f)\!:\!P$ for some $f \in O$, $P \subseteq D$. Non-factual values (no truth value) cannot be relations. Non-computable values (no effective method) cannot be computed by any system.
\end{proof}

\begin{corollary}[Algebraic Optimization]\label{cor:algebra}
Programs composed from FFP combining forms admit algebraic transformation via the laws of~\cite[Sec.\,12.2]{backus78}. Compilation from FFP objects to optimized evaluation preserves $\rho$-semantics.
\end{corollary}

\begin{proof}
The 24 laws of the FP algebra~\cite[Sec.\,12.2]{backus78} are equations between functional forms. An FFP object representing a program can be transformed to an equivalent object by applying these laws. The $\rho$-image is preserved because the laws hold for all objects in the domain. Ahead-of-time compilation from objects to evaluation trees is one such transformation.
\end{proof}

\section{Conclusion}

AREST composes FFP/AST~\cite{backus78} with REST~\cite{fielding00} via ORM\,2~\cite{halpin08}. Readings are FFP objects; $\rho$ maps them to executables; the population is Codd's named set~\cite{codd70} stored as Backus's FILE cell; transitions produce hypermedia. The eight results establish: the grammar is unambiguous (Theorem~\ref{thm:grammar}); specifications and executables are the same object (Theorem~\ref{thm:spec}); state transfer is complete in one application (Theorem~\ref{thm:complete}); HATEOAS is a $\theta_1$ projection (Theorem~\ref{thm:hateoas}); all domain values are derivable (Theorem~\ref{thm:derive}); all queries are expressible (Theorem~\ref{thm:adequate}); all derivations are verbalizable (Theorem~\ref{thm:deriv-complete}); and the scope is factual-and-computable, which is everything (Theorem~\ref{thm:scope}).

\begin{acks}
Claude (Anthropic) was used as a programming assistant during implementation and as a drafting aid during manuscript preparation. All architectural decisions and the composition of FFP/AST with REST were made by the human author.
\end{acks}

%%
%% Bibliography
%%
\bibliographystyle{ACM-Reference-Format}
\begin{thebibliography}{7}

\bibitem{backus78}
J. Backus.
\newblock Can Programming Be Liberated from the von Neumann Style? {A} Functional Style and Its Algebra of Programs.
\newblock \emph{Communications of the ACM}, 21(8):613--641, 1978.

\bibitem{church36}
A. Church.
\newblock An Unsolvable Problem of Elementary Number Theory.
\newblock \emph{American Journal of Mathematics}, 58(2):345--363, 1936.

\bibitem{codd70}
E.F. Codd.
\newblock A Relational Model of Data for Large Shared Data Banks.
\newblock \emph{Communications of the ACM}, 13(6):377--387, 1970.

\bibitem{fielding00}
R.T. Fielding.
\newblock \emph{Architectural Styles and the Design of Network-based Software Architectures}.
\newblock Doctoral dissertation, University of California, Irvine, 2000.

\bibitem{halpin08}
T. Halpin.
\newblock \emph{Information Modeling and Relational Databases}.
\newblock 2nd edition. Morgan Kaufmann, 2008.

\bibitem{halpin06}
T. Halpin and M. Curland.
\newblock Automated Verbalization for {ORM} 2.
\newblock In \emph{OTM Workshops}, pages 1181--1190, 2006.

\bibitem{halpin06b}
T. Halpin, M. Curland, and CS445 Class.
\newblock {ORM} 2 Constraint Verbalization: {P}art 1.
\newblock Technical Report ORM2-02, Neumont University, 2006.

\end{thebibliography}

\end{document}
